% Part: incompleteness
% Chapter: representability-in-q
% Section: introduction

\documentclass[../../../include/open-logic-section]{subfiles}

\begin{document}

\olfileid{inc}{req}{int}
\olsection{مقدمه}

قضایای ناتمامیت بر نظریه‌هایی اعمال می‌شوند که در آن‌ها بتوان
واقعیت‌های پایه دربارهٔ توابع محاسبه‌پذیر را بیان و اثبات کرد. نظریه‌ای
بسیار کمینه از این نوع را با نام «$\Th{Q}$» (یا گاهی «$Q$ رابینسون»،
به نام رافائل رابینسون) توصیف خواهیم کرد. خواهیم گفت
\emph{بازنمایی‌پذیر} بودن یک تابع در~$\Th{Q}$ به چه معناست و سپس
عبارت زیر را ثابت خواهیم کرد:
\begin{quote}
  یک تابع در~$\Th{Q}$ بازنمایی‌پذیر است اگر و تنها اگر
  محاسبه‌پذیر باشد.
\end{quote}
این از یک سو مدل دیگری از محاسبه‌پذیری به ما می‌دهد. اما از آن برای
نشان‌دادن تصمیم‌ناپذیربودن مجموعهٔ
$\Setabs{!A}{\Th{Q} \Proves !A}$ نیز استفاده خواهیم کرد، بدین صورت
که مسئلهٔ توقف را به آن کاهش می‌دهیم. تا پایان کار، نتایجی بسیار
قوی‌تر از این نیز ثابت کرده‌ایم.

زبان~$\Th{Q}$ زبان حساب است؛ $\Th{Q}$ از اصول موضوعهٔ زیر تشکیل
می‌شود (که باید همراه با دیگر اصول موضوعه و قواعد منطق مرتبه‌اول
دارای !!{identity} به کار روند):
\begin{align*}
& \lforall[x][\lforall[y][(\eq[x'][y'] \lif \eq[x][y])]] \tag{$!Q_1$}\\
& \lforall[x][\eq/[\Obj 0][x']] \tag{$!Q_2$}\\
& \lforall[x][(\eq[x][\Obj 0] \lor \lexists[y][\eq[x][y']])] \tag{$!Q_3$}\\
& \lforall[x][\eq[(x + \Obj 0)][x]] \tag{$!Q_4$}\\
& \lforall[x][\lforall[y][\eq[(x + y')][(x + y)']]] \tag{$!Q_5$}\\
& \lforall[x][\eq[(x \times \Obj 0)][\Obj 0]] \tag{$!Q_6$}\\
& \lforall[x][\lforall[y][\eq[(x \times y')][((x \times y) + x)]]] \tag{$!Q_7$}\\
& \lforall[x][\lforall[y][(x < y \liff \lexists[z][\eq[(z' + x)][y]])]] \tag{$!Q_8$}
\end{align*}
برای هر عدد طبیعی~$n$، عددنمای~$\num{n}$ را ترم
$\Obj{0}^{\prime\prime\ldots\prime}$ تعریف کنید که در مجموع~$n$
نشانهٔ پریم دارد. پس $\num{0}$ خودِ !!{constant}~$\Obj{0}$،
$\num{1}$ برابر~$\Obj{0}'$، $\num{2}$ برابر~$\Obj{0}''$ و به همین
ترتیب است.

$\Th{Q}$ به‌عنوان نظریه‌ای در حساب \emph{به‌شدت} ضعیف است؛ برای
نمونه، در آن حتی واقعیت‌های بسیار ساده‌ای چون
$\lforall[x][\eq/[x][x']]$ یا $\lforall[x][\lforall[y][(x+y)=(y+x)]]$
را نیز نمی‌توان ثابت کرد. اما خواهیم دید بخش بزرگی از علت جالب‌بودن
$\Th{Q}$ دقیقاً \emph{همین} ضعف آن است. در واقع، این نظریه فقط
به‌اندازه‌ای قوی است که قضیهٔ ناتمامیت بر آن اعمال شود. دلیل دیگر
اهمیت~$\Th{Q}$ آن است که مجموعهٔ اصول موضوعه‌اش \emph{متناهی} است.

نظریه‌ای قوی‌تر از~$\Th{Q}$، موسوم به \emph{حساب پئانو}~$\Th{PA}$،
با افزودن شِمای استقرا به~$\Th{Q}$ به دست می‌آید:
\[
(!A(\Obj 0) \land \lforall[x][(!A(x) \lif !A(x'))]) \lif \lforall[x][!A(x)]
\]
که در آن $!A(x)$ هر فرمولی است. اگر $!A(x)$ افزون بر~$x$ !!{variable}sی
آزاد دیگری داشته باشد، سورهای کلی را در آغاز می‌افزاییم تا همهٔ آن‌ها
بسته شوند (و نمونهٔ متناظر شِمای استقرا یک !!a{sentence} باشد). برای نمونه،
اگر $!A(x,y)$ !!{variable}~$y$ را نیز آزاد داشته باشد، نمونهٔ متناظر چنین
است:
\[
\lforall[y][((!A(\Obj 0) \land \lforall[x][(!A(x) \lif !A(x'))]) \lif
  \lforall[x][!A(x)])]
\]
با استفاده از نمونه‌های شِمای استقرا می‌توان از اصول موضوعهٔ
$\Th{PA}$ بسیار بیش از اصول موضوعهٔ~$\Th{Q}$ نتیجه گرفت. در واقع،
یافتن گزاره‌های «طبیعی» دربارهٔ اعداد طبیعی که در حساب پئانو اثبات‌پذیر
نباشند کار زیادی می‌طلبد!{}

\begin{defn}
\ollabel{defn:representable-fn}
  تابع $f(x_0,\ldots,x_k)$ از اعداد طبیعی به اعداد طبیعی را
  {\em بازنمایی‌پذیر در~$\Th{Q}$} می‌نامیم اگر فرمولی چون
  $!A_f(x_0,\dots,x_k,y)$ وجود داشته باشد به‌گونه‌ای که هرگاه
  $f(n_0,\dots,n_k)=m$، $\Th{Q}$ موارد زیر را ثابت کند:
\begin{enumerate}
\item\ollabel{defn:rep:a} $!A_f(\num{n_0}, \dots, \num{n_k}, \num{m})$
\item\ollabel{defn:rep:b} $\lforall[y][(!A_f(\num{n_0}, \dots,
\num{n_k}, y) \lif \num{m} = y)]$.
\end{enumerate}
\end{defn}

راه‌های دیگری نیز برای بیان تعریف هست؛ برای نمونه، می‌توانستیم به‌طور
هم‌ارز بخواهیم $\Th{Q}$ فرمول
$\lforall[y][(!A_f(\num{n_0},\dots,\num{n_k},y)
\liff\eq[y][\num{m}])]$ را ثابت کند.

\begin{thm}
\ollabel{thm:representable-iff-comp}
یک تابع در~$\Th{Q}$ بازنمایی‌پذیر است اگر و تنها اگر محاسبه‌پذیر باشد.
\end{thm}

اثبات قضیه دو جهت دارد. پس از آماده‌شدن حسابی‌سازی نحو، جهت چپ به
راست نسبتاً سرراست است. جهت دیگر کار بیشتری می‌طلبد. ایدهٔ پایه چنین
است: «بازگشتی عمومی» را معنای دقیق «محاسبه‌پذیر» می‌گیریم و نشان
می‌دهیم هر تابع بازگشتی عمومی در~$\Th{Q}$ بازنمایی‌پذیر است. به یاد
آورید تابعی بازگشتی عمومی است اگر بتوان آن را از~$\Zero$، تابع
جانشین~$\Succ$ و توابع افکنش~$\Proj{n}{i}$، با استفاده از ترکیب،
بازگشت اولیه و کمینه‌سازی منظم تعریف کرد. پس یک راه نشان‌دادن
بازنمایی‌پذیری همهٔ توابع بازگشتی عمومی در~$\Th{Q}$ آن است که نشان
دهیم توابع پایه بازنمایی‌پذیرند و هرگاه توابعی بازنمایی‌پذیر باشند،
توابعی که با ترکیب، بازگشت اولیه و کمینه‌سازی منظم از آن‌ها تعریف
می‌شوند نیز بازنمایی‌پذیرند. به بیان دیگر، می‌توان نشان داد توابع پایه
بازنمایی‌پذیرند و توابع بازنمایی‌پذیر نسبت به ترکیب، بازگشت اولیه و
کمینه‌سازی منظم «بسته‌اند». این تضمین می‌کند هر تابع بازگشتی عمومی
بازنمایی‌پذیر باشد.

معلوم می‌شود مرحله‌ای که باید در آن بسته‌بودن توابع بازنمایی‌پذیر نسبت
به بازگشت اولیه را نشان دهیم دشوار است. برای پرهیز از این مرحله، نخست
نشان می‌دهیم در واقع می‌توان بدون بازگشت اولیه کار کرد؛ یعنی هر تابع
بازگشتی عمومی را می‌توان تنها با ترکیب و کمینه‌سازی منظم از توابع پایه
تعریف کرد. بدین منظور نشان می‌دهیم بازگشت اولیه را می‌توان با یک
کمینه‌سازی منظمِ مشخص انجام داد. اما برای کارکردن این روش باید چند تابع
پایهٔ دیگر نیز بیفزاییم: جمع، ضرب و تابع مشخصهٔ رابطهٔ
این‌همانی~$\Char{=}$. سپس می‌توانیم قضیه را با نشان‌دادن این دو امر
ثابت کنیم: همهٔ \emph{این} توابع پایه در~$\Th{Q}$ بازنمایی‌پذیرند و
توابع بازنمایی‌پذیر نسبت به ترکیب و کمینه‌سازی منظم بسته‌اند.

\end{document}
